% Generated draft: requires mathematical and source review.
\documentclass[11pt]{article}
\usepackage[margin=25mm]{geometry}
\usepackage{fontspec}
\setmainfont{Latin Modern Roman}
\usepackage{amsmath,amssymb,mathtools,unicode-math}
\setmathfont{Latin Modern Math}
\usepackage{enumitem,needspace,xurl,hyperref}
\setlength{\parindent}{0pt}
\setlength{\parskip}{5pt}
\setlength{\emergencystretch}{3em}
\newcommand{\N}{\mathbb{N}}
\newcommand{\Z}{\mathbb{Z}}
\newcommand{\Q}{\mathbb{Q}}
\newcommand{\R}{\mathbb{R}}
\newcommand{\C}{\mathbb{C}}
\newcommand{\F}{\mathbb{F}}
\newcommand{\A}{\mathbb{A}}
\newcommand{\PP}{\mathbb{P}}
\newcommand{\E}{\mathbb{E}}
\newcommand{\eps}{\varepsilon}
\newcommand{\dd}{\,\mathrm{d}}
\begin{document}
\section*{UnsolvedMath ID: 30002017}
\textbf{Source problem code:} catalogue entry
\subsection{Definitions and mathematical statement}
Let \(G\) be a finite group. Write \(\operatorname{Irr}(G)\) for its ordinary irreducible complex characters and define
\[
\operatorname{cd}(G)=\{\chi(1):\chi\in\operatorname{Irr}(G)\}.
\]
This is a set of degrees, with repetitions discarded. A group is solvable if its derived series, defined by \(G^{(0)}=G\) and \(G^{(i+1)}=[G^{(i)},G^{(i)}]\), reaches the trivial group. The complex group algebra \(\C G\) consists of formal complex linear combinations of elements of \(G\), with multiplication extending the group multiplication. In the second question, isomorphisms of complex group algebras are interpreted as unital \(\C\)-algebra isomorphisms.

For a prime \(p\), a Sylow \(p\)-subgroup of \(G\) has order equal to the largest power of \(p\) dividing \(|G|\). If \(p\nmid |G|\), this subgroup is trivial and normal.

For the third question, use the conventional interpretation of a \(p\)-block: let \(k\) be an algebraically closed field of characteristic \(p\), and let \(B\) be a block of \(kG\), with primitive central block idempotent \(e_B\). Write \(\operatorname{Irr}(B)\) for the ordinary irreducible characters belonging to this block, using a splitting \(p\)-modular system to relate ordinary and modular characters. For a \(p\)-subgroup \(Q\leq G\), the Brauer homomorphism is
\[
\operatorname{Br}_Q:(kG)^Q\longrightarrow kC_G(Q),\qquad
\operatorname{Br}_Q\!\left(\sum_{g\in G}a_g g\right)=\sum_{g\in C_G(Q)}a_g g,
\]
where \((kG)^Q\) denotes the subalgebra fixed under conjugation by \(Q\). A \(B\)-Brauer pair is a pair \((Q,e_Q)\), where \(Q\) is a \(p\)-subgroup and \(e_Q\) is a block idempotent of \(kC_G(Q)\) satisfying \(\operatorname{Br}_Q(e_B)e_Q\ne0\). Put
\[
N_G(Q,e_Q)=\{g\in N_G(Q):g e_Qg^{-1}=e_Q\}.
\]
The block \(B\) is nilpotent if \(N_G(Q,e_Q)/C_G(Q)\) is a \(p\)-group for every \(B\)-Brauer pair. A defect group \(D\) of \(B\) is a \(p\)-subgroup maximal under inclusion subject to \(\operatorname{Br}_D(e_B)\ne0\). Write \(|G|=p^a m\), with \(p\nmid m\), and \(|D|=p^d\). For \(\chi\in\operatorname{Irr}(B)\), its height \(h_B(\chi)\) is defined by
\[
v_p(\chi(1))=a-d+h_B(\chi),
\]
where \(v_p(n)\) is the exponent of \(p\) in the positive integer \(n\). Set \(\operatorname{Irr}_0(B)=\{\chi\in\operatorname{Irr}(B):h_B(\chi)=0\}\).

The supplied record asks whether each of the following assertions holds:
\begin{enumerate}
\item For all finite groups \(G,H\),
\[
\operatorname{cd}(G)=\operatorname{cd}(H)
\quad\Longrightarrow\quad
\bigl(G\text{ is solvable}\iff H\text{ is solvable}\bigr).
\]
\item For every prime \(p\) and all finite groups \(G,H\),
\[
\C G\cong_{\C\text{-alg}}\C H
\quad\Longrightarrow\quad
\left(
\begin{array}{c}
G\text{ has a normal Sylow }p\text{-subgroup}
\end{array}
\iff
\begin{array}{c}
H\text{ has a normal Sylow }p\text{-subgroup}
\end{array}
\right).
\]
The same prime \(p\) is used on both sides.
\item For every prime \(p\), every finite group \(G\), and every \(p\)-block \(B\) of \(G\),
\[
B\text{ is nilpotent}
\quad\Longleftrightarrow\quad
\chi(1)=\psi(1)\quad\text{for all }\chi,\psi\in\operatorname{Irr}_0(B).
\]
Equality here means equality of the full ordinary character degrees.
\end{enumerate}
The block-theoretic setting and algebra-isomorphism convention above make explicit standard readings that the supplied statement does not spell out. The record supplies no additional restrictions on groups, primes, or blocks. Its background discusses Dade's Projective Conjecture without specifying how that discussion relates to these three questions; that unresolved source discrepancy is not an additional conjecture asserted here.
\subsection{Short English statement}
The record asks three questions: whether the set of ordinary irreducible character degrees determines solvability, whether a complex group algebra determines the existence of a normal Sylow \(p\)-subgroup for each fixed prime \(p\), and whether a block is nilpotent exactly when all its height-zero ordinary characters have equal degrees. The supplied open label and literature assessment require independent review, especially because the background instead discusses Dade's Projective Conjecture.
\subsection{Sources}
\begin{itemize}
\item Local UnsolvedMath registry, supplied complete record with permanent ID \(30002017\), title \emph{Character Degrees, Group Algebras, and Nilpotent Blocks}; exact source problem code \texttt{catalogue entry}. Collection: Oberwolfach Reports Open Problems. The record gives the proposed year as 2012 and names no individual proposer.
\item Supplied report citation: \emph{Representations of Finite Groups} (2012), Oberwolfach Reports, pp.~963--1019, DOI \texttt{10.4171/owr/2012/16}. \url{https://doi.org/10.4171/owr/2012/16}. This is supplied report-level provenance; no precise page or passage for these three questions is provided.
\item Supplied literature source: \emph{Reduction theorem for Dade's projective conjecture}. \url{https://ems.press/journals/jems/articles/14667}.
\item Supplied literature source: \emph{Character Triple reduction}. \url{https://arxiv.org/abs/2402.10632}.
\item Supplied literature check date: \texttt{2026-08-21}. The supplied dated curation labels the problem open and assesses that the fragment invokes Dade's Projective Conjecture, that reductions and broad special cases are known, and that the conjecture is not proved for all finite groups. This is supplied dated curation, not an independently verified assessment. The supplied background expressly requests independent verification of time-sensitive status and cited claims.
\item The registry labels the statement \texttt{corrected\_verified} and reports a manual semantic, typographic, and LaTeX review claiming that its cleanup is self-contained and meaning-preserving. These are supplied registry assertions, not mathematical or source verification performed for this draft.
\end{itemize}
\end{document}
